\documentclass[12pt,a4paper]{article}
\usepackage[UTF8]{ctex}
\usepackage{amsmath,amssymb,amsthm}
\usepackage{geometry}
\usepackage{graphicx}
\usepackage{hyperref}
\usepackage{bm}

\geometry{margin=2.5cm}

\title{空间惯性矩阵$G$的定义与推导}
\author{Modern Robotics}
\date{}

\begin{document}

\maketitle

\section{引言}

空间惯性矩阵$G$（spatial inertia matrix）是描述刚体惯性特性的$6 \times 6$矩阵，它将twist（速度）与wrench（力）联系起来。本文档详细说明$G$的定义、推导过程及其物理意义。

\section{基本定义}

\subsection{空间惯性矩阵的定义}

对于刚体，其空间惯性矩阵$G \in \mathbb{R}^{6 \times 6}$定义为：

\begin{equation}
G = \begin{bmatrix}
I & 0 \\
0 & mI
\end{bmatrix}
\label{eq:G_definition}
\end{equation}

其中：
\begin{itemize}
\item $I \in \mathbb{R}^{3 \times 3}$：刚体在质心处的转动惯量矩阵（rotational inertia matrix）
\item $m \in \mathbb{R}$：刚体的质量
\item $I$（右下角）：$3 \times 3$单位矩阵
\end{itemize}

\textbf{注意}：这个定义假设坐标系固定在刚体的质心处，且坐标轴与主惯性轴对齐（或至少$I$是已知的）。

\section{从单刚体动力学推导}

\subsection{问题设置}

考虑一个由多个点质量组成的刚体。设：
\begin{itemize}
\item 点质量$i$的质量为$m_i$
\item 总质量：$m = \sum_i m_i$
\item 点质量$i$在体坐标系$\{b\}$中的位置为$\bm{r}_i = (x_i, y_i, z_i)^T$
\item 体坐标系$\{b\}$的原点位于质心，即：$\sum_i m_i \bm{r}_i = 0$
\end{itemize}

\subsection{刚体的运动}

设刚体以twist $V_b = (\omega_b, v_b)$运动，其中：
\begin{itemize}
\item $\omega_b \in \mathbb{R}^3$：角速度
\item $v_b \in \mathbb{R}^3$：线速度（质心的速度）
\end{itemize}

点质量$i$的速度为：
\begin{equation}
\dot{\bm{p}}_i = v_b + \omega_b \times \bm{p}_i
\end{equation}

其中$\bm{p}_i(t)$是点质量$i$在惯性系中的位置。

\subsection{点质量的加速度}

对速度求导：
\begin{align}
\ddot{\bm{p}}_i &= \frac{d}{dt}[v_b + \omega_b \times \bm{p}_i] \\
&= \dot{v}_b + \dot{\omega}_b \times \bm{p}_i + \omega_b \times \dot{\bm{p}}_i \\
&= \dot{v}_b + \dot{\omega}_b \times \bm{p}_i + \omega_b \times (v_b + \omega_b \times \bm{p}_i) \\
&= \dot{v}_b + \dot{\omega}_b \times \bm{r}_i + [\omega_b]v_b + [\omega_b]^2 \bm{r}_i
\end{align}

使用反对称矩阵表示：
\begin{equation}
\ddot{\bm{p}}_i = \dot{v}_b + [\dot{\omega}_b]\bm{r}_i + [\omega_b]v_b + [\omega_b]^2 \bm{r}_i
\label{eq:point_acceleration}
\end{equation}

\subsection{作用在点质量上的力}

根据牛顿第二定律，作用在点质量$i$上的力为：
\begin{equation}
\bm{f}_i = m_i \ddot{\bm{p}}_i = m_i(\dot{v}_b + [\dot{\omega}_b]\bm{r}_i + [\omega_b]v_b + [\omega_b]^2 \bm{r}_i)
\end{equation}

\subsection{总力和总力矩}

作用在刚体上的总力为：
\begin{equation}
\bm{f}_b = \sum_i \bm{f}_i = \sum_i m_i \ddot{\bm{p}}_i
\end{equation}

作用在刚体上的总力矩（关于质心）为：
\begin{equation}
\bm{m}_b = \sum_i \bm{r}_i \times \bm{f}_i = \sum_i [\bm{r}_i] \bm{f}_i
\end{equation}

\subsection{线动力学的推导}

计算总力：
\begin{align}
\bm{f}_b &= \sum_i m_i(\dot{v}_b + [\dot{\omega}_b]\bm{r}_i + [\omega_b]v_b + [\omega_b]^2 \bm{r}_i) \\
&= \sum_i m_i \dot{v}_b + \sum_i m_i [\dot{\omega}_b]\bm{r}_i + \sum_i m_i [\omega_b]v_b + \sum_i m_i [\omega_b]^2 \bm{r}_i
\end{align}

利用质心性质：$\sum_i m_i \bm{r}_i = 0$，我们有：
\begin{itemize}
\item $\sum_i m_i [\dot{\omega}_b]\bm{r}_i = [\dot{\omega}_b] \sum_i m_i \bm{r}_i = 0$
\item $\sum_i m_i [\omega_b]^2 \bm{r}_i = [\omega_b]^2 \sum_i m_i \bm{r}_i = 0$
\end{itemize}

因此：
\begin{equation}
\bm{f}_b = m \dot{v}_b + m [\omega_b] v_b = m(\dot{v}_b + [\omega_b] v_b)
\label{eq:linear_dynamics}
\end{equation}

\subsection{转动力学的推导}

计算总力矩：
\begin{align}
\bm{m}_b &= \sum_i [\bm{r}_i] m_i(\dot{v}_b + [\dot{\omega}_b]\bm{r}_i + [\omega_b]v_b + [\omega_b]^2 \bm{r}_i) \\
&= \sum_i m_i [\bm{r}_i]\dot{v}_b + \sum_i m_i [\bm{r}_i][\dot{\omega}_b]\bm{r}_i + \sum_i m_i [\bm{r}_i][\omega_b]v_b + \sum_i m_i [\bm{r}_i][\omega_b]^2 \bm{r}_i
\end{align}

利用质心性质：
\begin{itemize}
\item $\sum_i m_i [\bm{r}_i]\dot{v}_b = [\sum_i m_i \bm{r}_i]\dot{v}_b = 0$
\item $\sum_i m_i [\bm{r}_i][\omega_b]v_b = [\sum_i m_i \bm{r}_i][\omega_b]v_b = 0$
\end{itemize}

对于第二项，使用$[\bm{r}_i][\dot{\omega}_b]\bm{r}_i = -[\bm{r}_i]^2 \dot{\omega}_b$：
\begin{equation}
\sum_i m_i [\bm{r}_i][\dot{\omega}_b]\bm{r}_i = -\sum_i m_i [\bm{r}_i]^2 \dot{\omega}_b
\end{equation}

对于第四项，使用$[\bm{r}_i][\omega_b]^2 \bm{r}_i = -[\bm{r}_i]^T [\omega_b]^T [\bm{r}_i] \omega_b$，经过整理：
\begin{equation}
\sum_i m_i [\bm{r}_i][\omega_b]^2 \bm{r}_i = -[\omega_b] \sum_i m_i [\bm{r}_i]^2 \omega_b
\end{equation}

定义转动惯量矩阵：
\begin{equation}
I_b = -\sum_i m_i [\bm{r}_i]^2
\label{eq:inertia_definition}
\end{equation}

则：
\begin{equation}
\bm{m}_b = I_b \dot{\omega}_b + [\omega_b] I_b \omega_b
\label{eq:rotational_dynamics}
\end{equation}

\subsection{转动惯量矩阵的展开}

展开$[\bm{r}_i]^2$：
\begin{equation}
[\bm{r}_i]^2 = \begin{bmatrix}
0 & -z_i & y_i \\
z_i & 0 & -x_i \\
-y_i & x_i & 0
\end{bmatrix}^2 = \begin{bmatrix}
-(y_i^2 + z_i^2) & x_i y_i & x_i z_i \\
x_i y_i & -(x_i^2 + z_i^2) & y_i z_i \\
x_i z_i & y_i z_i & -(x_i^2 + y_i^2)
\end{bmatrix}
\end{equation}

因此：
\begin{equation}
I_b = -\sum_i m_i [\bm{r}_i]^2 = \begin{bmatrix}
\sum_i m_i(y_i^2 + z_i^2) & -\sum_i m_i x_i y_i & -\sum_i m_i x_i z_i \\
-\sum_i m_i x_i y_i & \sum_i m_i(x_i^2 + z_i^2) & -\sum_i m_i y_i z_i \\
-\sum_i m_i x_i z_i & -\sum_i m_i y_i z_i & \sum_i m_i(x_i^2 + y_i^2)
\end{bmatrix}
\end{equation}

写成标准形式：
\begin{equation}
I_b = \begin{bmatrix}
I_{xx} & I_{xy} & I_{xz} \\
I_{xy} & I_{yy} & I_{yz} \\
I_{xz} & I_{yz} & I_{zz}
\end{bmatrix}
\end{equation}

其中：
\begin{align}
I_{xx} &= \sum_i m_i(y_i^2 + z_i^2) \\
I_{yy} &= \sum_i m_i(x_i^2 + z_i^2) \\
I_{zz} &= \sum_i m_i(x_i^2 + y_i^2) \\
I_{xy} &= -\sum_i m_i x_i y_i \\
I_{xz} &= -\sum_i m_i x_i z_i \\
I_{yz} &= -\sum_i m_i y_i z_i
\end{align}

\subsection{组合成空间形式}

将线动力学(\ref{eq:linear_dynamics})和转动力学(\ref{eq:rotational_dynamics})组合：

\begin{equation}
\begin{bmatrix}
\bm{m}_b \\
\bm{f}_b
\end{bmatrix} = \begin{bmatrix}
I_b \dot{\omega}_b + [\omega_b] I_b \omega_b \\
m \dot{v}_b + m [\omega_b] v_b
\end{bmatrix}
\end{equation}

写成矩阵形式：
\begin{equation}
\begin{bmatrix}
\bm{m}_b \\
\bm{f}_b
\end{bmatrix} = \begin{bmatrix}
I_b & 0 \\
0 & mI
\end{bmatrix} \begin{bmatrix}
\dot{\omega}_b \\
\dot{v}_b
\end{bmatrix} - \begin{bmatrix}
-[\omega_b] I_b \omega_b \\
-m [\omega_b] v_b
\end{bmatrix}
\end{equation}

使用$[\omega_b]^T = -[\omega_b]$，整理得：
\begin{equation}
\begin{bmatrix}
\bm{m}_b \\
\bm{f}_b
\end{bmatrix} = \begin{bmatrix}
I_b & 0 \\
0 & mI
\end{bmatrix} \begin{bmatrix}
\dot{\omega}_b \\
\dot{v}_b
\end{bmatrix} - \begin{bmatrix}
[\omega_b] & 0 \\
[v_b] & [\omega_b]
\end{bmatrix}^T \begin{bmatrix}
I_b \omega_b \\
m v_b
\end{bmatrix}
\end{equation}

定义空间惯性矩阵：
\begin{equation}
G_b = \begin{bmatrix}
I_b & 0 \\
0 & mI
\end{bmatrix}
\label{eq:G_final}
\end{equation}

则动力学方程可以写成：
\begin{equation}
F_b = G_b \dot{V}_b - [\text{ad}_{V_b}]^T G_b V_b
\label{eq:spatial_dynamics}
\end{equation}

其中$F_b = (\bm{m}_b, \bm{f}_b)$是wrench，$V_b = (\omega_b, v_b)$是twist。

\section{空间惯性矩阵的性质}

\subsection{对称性和正定性}

\begin{enumerate}
\item \textbf{对称性}：$G_b$是对称矩阵，因为$I_b$是对称的
\item \textbf{正定性}：$G_b$是正定矩阵，因为：
\begin{itemize}
\item $I_b$是正定的（转动惯量矩阵）
\item $m > 0$（质量为正）
\item 因此对于任意非零twist $V_b$，有$V_b^T G_b V_b > 0$
\end{itemize}
\end{enumerate}

\subsection{与动能的关系}

刚体的动能为：
\begin{equation}
K = \frac{1}{2} \sum_i m_i \|\dot{\bm{p}}_i\|^2
\end{equation}

可以证明：
\begin{equation}
K = \frac{1}{2} \omega_b^T I_b \omega_b + \frac{1}{2} m v_b^T v_b = \frac{1}{2} V_b^T G_b V_b
\label{eq:kinetic_energy}
\end{equation}

这证明了空间惯性矩阵$G_b$完全描述了刚体的惯性特性。

\subsection{与动量的关系}

刚体的空间动量为：
\begin{equation}
P_b = G_b V_b = \begin{bmatrix}
I_b \omega_b \\
m v_b
\end{bmatrix}
\label{eq:spatial_momentum}
\end{equation}

其中：
\begin{itemize}
\item $I_b \omega_b$：角动量
\item $m v_b$：线动量
\end{itemize}

\section{在不同坐标系中的表示}

\subsection{坐标系变换}

如果空间惯性矩阵$G_b$在坐标系$\{b\}$（质心）中定义，那么在另一个坐标系$\{a\}$中的表示为：

\begin{equation}
G_a = [\text{Ad}_{T_{ba}}]^T G_b [\text{Ad}_{T_{ba}}]
\label{eq:G_transform}
\end{equation}

其中$T_{ba}$是从坐标系$\{b\}$到$\{a\}$的变换。

\textbf{推导}：

由于动能是标量，不依赖于坐标系：
\begin{equation}
\frac{1}{2} V_a^T G_a V_a = \frac{1}{2} V_b^T G_b V_b
\end{equation}

使用伴随变换：$V_b = [\text{Ad}_{T_{ba}}] V_a$，因此：
\begin{align}
\frac{1}{2} V_a^T G_a V_a &= \frac{1}{2} ([\text{Ad}_{T_{ba}}] V_a)^T G_b ([\text{Ad}_{T_{ba}}] V_a) \\
&= \frac{1}{2} V_a^T [\text{Ad}_{T_{ba}}]^T G_b [\text{Ad}_{T_{ba}}] V_a
\end{align}

因此：
\begin{equation}
G_a = [\text{Ad}_{T_{ba}}]^T G_b [\text{Ad}_{T_{ba}}]
\end{equation}

\subsection{Steiner定理的推广}

如果坐标系$\{a\}$不在质心，而是在距离质心$\bm{p}$的位置，则：

\begin{equation}
G_a = \begin{bmatrix}
I_b + m([\bm{p}]^T [\bm{p}] - \|\bm{p}\|^2 I) & m [\bm{p}]^T \\
m [\bm{p}] & mI
\end{bmatrix}
\end{equation}

这可以简化为：
\begin{equation}
G_a = [\text{Ad}_{T_{ba}}]^T G_b [\text{Ad}_{T_{ba}}]
\end{equation}

其中$T_{ba}$是从质心到$\{a\}$的变换。

\section{在机器人动力学中的应用}

\subsection{连杆的空间惯性矩阵}

对于机器人连杆$i$，其空间惯性矩阵$G_i$定义为：

\begin{equation}
G_i = \begin{bmatrix}
I_i & 0 \\
0 & m_i I
\end{bmatrix}
\label{eq:link_inertia}
\end{equation}

其中：
\begin{itemize}
\item $I_i$：连杆$i$在质心处的转动惯量矩阵（在坐标系$\{i\}$中表示）
\item $m_i$：连杆$i$的质量
\item 坐标系$\{i\}$固定在连杆$i$的质心处
\end{itemize}

\subsection{在牛顿-欧拉算法中的使用}

在牛顿-欧拉逆向动力学算法中，$G_i$用于计算连杆$i$的动力学：

\begin{equation}
F_i = G_i \dot{V}_i - [\text{ad}_{V_i}]^T G_i V_i
\end{equation}

其中：
\begin{itemize}
\item $F_i$：作用在连杆$i$上的wrench
\item $V_i$：连杆$i$的twist
\item $\dot{V}_i$：连杆$i$的twist的导数
\end{itemize}

\section{总结}

\begin{enumerate}
\item \textbf{定义}：$G = \begin{bmatrix} I & 0 \\ 0 & mI \end{bmatrix}$，其中$I$是转动惯量矩阵，$m$是质量

\item \textbf{物理意义}：
\begin{itemize}
\item 描述刚体的惯性特性
\item 将twist与wrench联系起来
\item 与动能和动量直接相关
\end{itemize}

\item \textbf{性质}：
\begin{itemize}
\item 对称正定矩阵
\item 动能：$K = \frac{1}{2} V^T G V$
\item 动量：$P = G V$
\end{itemize}

\item \textbf{坐标系变换}：$G_a = [\text{Ad}_{T_{ba}}]^T G_b [\text{Ad}_{T_{ba}}]$

\item \textbf{应用}：在机器人动力学中用于计算每个连杆的惯性力
\end{enumerate}
\end{document}